\scanpage{114}
Par
\[
\operatorname{Cons}(X,E)(Y)=\operatorname{App}_{\mathrm{cons}}(Y,(X,E))
\]
\deleted{\unsure{On a}}

Si on a $(X,E)$ et $(X',E')$, on a.
\[
\operatorname{App}_{\mathrm{cons}}((X,E),(X',E'))\xrightarrow{\sim}\operatorname{Hom}(\operatorname{Cons}(X,E),\operatorname{Cons}(X',E'))
\]
On obtient ainsi des foncteurs pleinement fidèles

\diagram{\begin{tikzcd}[ampersand replacement=\&,column sep=large,row sep=huge]
(\mathrm{Cons\ Aff})
 \arrow[r]
 \arrow[r,bend left=42,"{\substack{\text{\unsure{épimorphismes|morphismes les}}\\\text{\unsure{cons. définies par q. cpts}...}}}"]
 \arrow[rr,bend right=35,"{\text{canonique}}"']
 \arrow[rrr,bend right=58,"{\text{can}}"']
\& \begin{gathered}(\text{Couples }(X,E))\\\text{$E$ loc. constr.}\\\text{dans $\unsure{X|E}$ quasép}\end{gathered}
 \arrow[r]
\& \underbrace{\operatorname{Ind}(\mathrm{Cons\ Aff})'}_{\substack{\text{systèmes}\\\text{inductifs}\\\text{\unsure{immersifs}---}}}
 \arrow[r]
 \arrow[r,bend right=40,"{\text{can}}"']
\& \underline{\operatorname{Hom}}((\mathrm{Cons\ Aff})^\circ,\unsure{\mathrm{Ens}})
\end{tikzcd}}
\margin{113}
